\documentclass[12pt]{article}

\usepackage[a4paper, total={6in, 8in}]{geometry}
\usepackage{textcomp}

\begin{document}
\setlength{\parindent}{0pt}

\def \t {\textrightarrow}
\def \v {\vspace{1em}}
\def \bi {\begin{itemize}}
\def \ei {\end{itemize}}
\def \s[#1] {\section*{#1}}
\def \ss[#1] {\subsection*{#1}}
\def \sss[#1] {\subsubsection*{#1}}

\s[Ungaretti]

\ss[Vita]
Nasce nel ? e muore nel ? \\
È un poeta che vive in diversi luoghi \t la sua esperienze si dirama dal Brasile, poi a Parigi conosce le avanguardie \t e sperimenta la passione per gli abissi \\
Il suo periodo in trincea \t porta alla questione dell'interventismo oppure no \t sperimenta quanto il dolore sia pregnante, e vive sul campo la sofferenza della guerra \\
Bisogna anche ricordare alcuni fatti della sua biografia \t suo figlio antonietto muore \t le sue raccolte "Allegria di naufragi" e "Porto sepolto" \\
Anche in carducci questo dolore è presente \t non si puo stabilire i parametri della gestione del dolore \t quindi non si possono confrontare i due dolori \\
In Carducci non c'è il conforto della fede, mentre in Ungaretti si, e offre un margine di accettazione che pero non si esaurisce \\
L'importanza della sua frequentazione parigina si vede anche in una influenza del futurismo, che viene dalla letteratura (Oliner, Malarme) ma anche dall arte (Picasso, De Chirico, Renoir) \t molte sono le provocazioni artistiche che gli arrivano

\v

Quando il pieno della guerra lo coinvolge, la sua prospettiva di vita cambia \\
C'è una differenza tra l ungarettima prima e dopo la guerra \t ha una cattedra a San Paolo, mentre negli anni 70 ottiene riconoscimento a Milano, e ottiene cattedra poi alla Sapienza di Roma \\
La sua parabola si configura nella raccolta "Dolore" \t sofferenza profonda radicata negli abissi \\
La sua speranza è quella di una terra promessa \t ma una terra dove l uomo possa fare ritorno e possa vedere il compimento di un cerchio \\
Questa parabola comporta la parola che diventa grido strozzato, e il grande lavoro filologico che l'autore compie su se stesso (in "Vita di un uomo") \\
"Allegria di naufragi" = sorta di ossimoro, come si puo trovare l'allegria in un evento terribile? \t in questo titolo c'è tutto il dramma della devastazione del primo 900 \t il suo senso di smarrimento è quello dell uomo che vede la guerra e si colloca nel disfacimento di tutto \\
"Porto sepolto" = figura del poeta che è palombaro, e titolo ricorda le civiltà sommerse \t e quindi il far riaffiorare le verità, che a un non so che di escatologico e rivelatore \t con provenienza socratica e riporta alla psicanalisi, che fa riemerge gli irrisolti \\
Psicanalisi pero da lettura al fenomeno, mentra qua si da chiarezza sul fenomeno \t si ha la pretesa di portare a galla il dolore, non di esorcizzarlo come nella pscianalisi \\
"Vita di un uomo" = tutto cio che leggiamo si trova qua, dove celebra l'atto finale

\s[I testi della guerra]
\ss[Veglia]
Scrittura che si fa illuminazione, equilibrio spazio bianco, poca punteggiatura, e l'urgenza di esprimersi \\
Scrive in guerra su fogli che trova \t a notte inoltrata o al mattimo presto \t scrive del particolare conferendo ad esso messaggio universale (che non si legge + in montale, che è emeretismo esistenziale) \\
Qua a differenza di montale la valenza metafisica è ancora presente \t e aleggia il sentimento del tempo e del senso religioso, che pero non coincide all adesione a una tale religione \t ma individua le domande di senso \\
Tutti noi abbiamo un senso religioso, che non è dio, ma è "perchè sono qui", "chi sono" \t che permettono di vivere e non di vegetare \\
Ogni componimento è caratterizzato dal luogo e data di composizione \t "Veglia" analizzato di Contini corrisponde all abbandono del ? dannunziamo, composta a Quarto nel 19??

\v

Si osserva subito lo spazio \t e che i versi sono caratterizzati dall'assenza di punteggiatura \\
L'architettura interna del testo si ridisegna in una mappa rinnovata (equilibrio bianco scritta) \t il verso 6 abbiamo verso con una sola parola per esempio, anche al 10 e al 15 \\
Al verso 4 anche, con solo massacrato \\
Preponderanza di participi perfetti \t sottolinea l'impotenza del poeta di fronte al massacro \\
Buttato \t sottolinea una non volonta ma una circostanza in cui ci si trova \\
Plenilunio \t ricorda leopardi, ma non è \t è una luna che assiste alla fine di una vita \t e un dolore che penetra nelle viscere, e nella bocca digrignata di un compagno di guerra \\
Bocca digrignata è espressione palpabile del dolore \t e questo dolore diventa del poeta: "penetrata del mio silenzio" \t quel vissuto condiviso diventa vissuto anche dell'autore \\
Sue mani - mio silenzio (congestionate dal freddo) \t questo dolore attraverso la congestione \t le mani sono artefici di tanto che l animo puo fare, mentre qua sono private di ogni potenzialita \t è la fine del tutto \\
In questa totale assenza di speranza e in questa irremidibilità del tutto (la morte), in chiave personale l'amore si leva \\
L'amore nasce dalla radice di dolore penetrato nel suo animo \t e nell esperienza di morte, lui assapora il senso dell'esistere \t non sono mai stato tanto attaccato alla vita  \\
Qua la guerra viene vissuta come universale ed educativa nei confronti del senso della vita \t che emerge dal dolore e dalla resilienza di fronte al dolore \\
Lui non si guarda vivere \t il suo fulcro non è guardare fuori da se, perche quel dolore è diventato parte di se \t immagine totale dell'empatia, e di una guerra che viene vissuta dall'autore 

\ss[Fratelli]
Testo breve, fulmineita della parola e frammento lirico \t ritmo spezzato da architettura testuale, che privilegia le pause (spazi bianchi) e evidenzia la parola fratelli \\
Il combattere è sentito \t non è quello del soldato mercenario, il combattere è espressione di una fratellanza \t dove l uomo si trova parte di un disegno piu grande \\
La fratellanza azzera la provenienza, e anche questo è stato composto 15 luglio del 1916 \t le date contestualizzano pregnanza di una esperienza \\
Lui sceglie di mettere queste date, nel rispetto della testualita e nell'aspetto filologico \t esprime la ricorrenza che ha un senso \\
Secondo una visione hegeliana della storia, non è il dolore indivduale, ma è il dolore di tutti \\
L'uomo è dolore, leopardi infatti afferma che l'uomo nasce a fatica \\
È proprio dei grandi attraversare il dolore ed esorcizzarlo

\v

Versi 2, 4, e gli ultimi due versi (fragilità è isolato insieme a fratelli) \\
Inizia con una domanda (sottointende che senso ha saperlo) \\
Dire la parola fratelli significa pronunciare una parola che si teme \t aria descritta come presa in spasmi \t che si ricollegano anche al pianto \\
Verso numero 6 \t sinestesia e (non prosopopea, aria non parla) \\
L uomo fa esperienza del suo essere nulla, nel momento in cui puo solo afferrare un tutto con gli altri \t una chiave diversa si vede nella "Ginestra" di Leopardi \\
L'uomo si mostra come non nato per il combattere \\


\end{document}
